\chapter{関係データベースの設計 -- 正規化}

「IT Text データベースの基礎」（吉川正俊著）\cite{吉川本}で述べられているように，理想的なデータベースは以下の性質を有するデータベーススキーマを持ちます：
\begin{enumerate}
  \item データの冗長性がない
  \item データの一貫性制約の保持が容易である
\end{enumerate}
関係データベースにおいては，上記(1)(2)の要件を満たすように関係スキーマを定義することを\strong{正規化（normalization）}と呼びます．
望ましい関係スキーマ，つまり正規化された関係スキーマが設計できれば，関係データベースから\strong{冗長性が排除され（データの重複がなく），データを正しく矛盾なく管理・処理}できるようになります．

前章では実体関連モデルを用いた関係スキーマの設計方法について説明しましたが，実はこの方法では十分な正規化が保証されません．
実体関連モデルの変換のみで適切な関係スキーマを得るためには，実体関連図の設計段階で正規化を意識した設計を行う必要があり，設計の質は設計者の経験やスキルに依存します．

このような問題を解決するのが，正規化の理論に基づいた関係スキーマの設計方法です．
本章では，実体関連モデルの変換による関係スキーマを補完する，データ従属性に基づく正規化手法について説明します．
